\chapter[Lezione 7 --- 14 ottobre 2026]{Lezione 7: numerabilità, operazioni sulle funzioni e punti di accumulazione}
\label{chap:lezione-07}

\section{Numerabilità degli insiemi}
\label{sec:lez07-numerabilita}

La nozione di applicazione biiettiva, introdotta nella
\cref{def:lez06-biettiva}, permette di confrontare insiemi
anche quando contengono infiniti elementi. Due insiemi
hanno la stessa \emph{cardinalità} se esiste una biiezione
tra essi.

\begin{definition}[Insieme numerabile]
\label{def:lez07-numerabilita}
Un insieme \(A\) si dice \emph{numerabile} se esiste
un'applicazione biiettiva
\[
  \varphi\colon\N\longrightarrow A.
\]
Un insieme infinito che non ammette una tale biiezione
si dice \emph{non numerabile}.
\end{definition}

\begin{proposition}[Numerabilità dei naturali pari]
\label{prop:lez07-pari}
L'insieme dei numeri naturali pari
\[
  P:=\{2m:m\in\N\}
\]
è numerabile. In particolare, l'applicazione
\[
  \varphi\colon\N\longrightarrow P,
  \qquad m\longmapsto 2m,
\]
è biiettiva.
\end{proposition}

\begin{proof}
Si verificano separatamente la suriettività e
l'iniettività dell'applicazione.

\emph{Suriettività.}
Sia \(p\in P\) arbitrario. Per definizione di \(P\),
esiste \(m\in\N\) tale che
\[
  p=2m=\varphi(m).
\]
Ogni elemento di \(P\) possiede quindi almeno un
antecedente in \(\N\). Per la
\cref{def:lez06-suriettiva}, \(\varphi\) è suriettiva.

\emph{Iniettività.}
Siano \(m_1,m_2\in\N\) con \(m_1\ne m_2\).
Supponiamo per assurdo che
\(\varphi(m_1)=\varphi(m_2)\). Allora
\[
  2m_1=2m_2
  \quad\Longrightarrow\quad
  m_1=m_2,
\]
in contraddizione con l'ipotesi \(m_1\ne m_2\).
Pertanto \(\varphi(m_1)\ne\varphi(m_2)\),
quindi \(\varphi\) è iniettiva.

Essendo sia suriettiva sia iniettiva,
\(\varphi\) è biiettiva.
\end{proof}

Analogamente, l'insieme dei naturali dispari
\[
  D:=\{2m+1:m\in\N\}
\]
è numerabile mediante l'applicazione
\[
  \psi\colon\N\longrightarrow D,
  \qquad m\longmapsto 2m+1.
\]
La suriettività segue dalla definizione di \(D\);
l'iniettività dall'implicazione
\[
  2m_1+1=2m_2+1
  \quad\Longrightarrow\quad
  m_1=m_2.
\]

Poiché \(P\subsetneq\N\) (infatti \(1\notin P\)),
la \cref{prop:lez07-pari} mostra che può esistere
una biiezione tra un insieme infinito e un suo
sottoinsieme proprio.

\begin{proposition}[Numerabilità degli insiemi numerici]
\label{prop:lez07-insiemi-numerici}
Gli insiemi \(\N\), \(\Z\) e \(\Q\) sono numerabili,
mentre \(\R\) non è numerabile.
\end{proposition}

Per \(\N\) è sufficiente considerare l'applicazione
identità. La numerabilità di \(\Z\) e \(\Q\) e la
non numerabilità di \(\R\) sono qui enunciate
senza dimostrazione. Quest'ultimo risultato
è dovuto a Cantor.

La numerabilità non va confusa con la densità:
in particolare, \(\Q\) è numerabile pur essendo
denso in \(\R\), come stabilito nel
\cref{thm:densita-razionali}.

\begin{example}[Albergo di Hilbert]
\label{ex:lez07-hilbert}
Si consideri un albergo con infinite camere,
numerate mediante i numeri naturali, tutte occupate.

Sebbene non vi siano camere libere, è possibile
accogliere un nuovo ospite trasferendo
simultaneamente ciascun occupante dalla camera
\(n\) alla camera \(n+1\), per ogni \(n\in\N\).

L'assegnazione
\[
  n\longmapsto n+1
\]
è iniettiva, quindi nessuna camera riceve più
di un ospite. Dopo lo spostamento, la camera
\(0\) rimane libera e può accogliere il nuovo
arrivato.

L'assegnazione costituisce una biiezione tra
\(\N\) e \(\N\setminus\{0\}\).

\begin{figure}[H]
  \centering
  \begin{tikzpicture}[unipd/diagram,x=1cm,y=1cm]
  \node[left] at (-.6,0) {Prima};
  \node[left] at (-.6,-1.65) {Dopo};
  \foreach \n in {0,...,5} {
    \draw[unipd-primary] ({\n-.38},-.35) rectangle ({\n+.38},.35);
    \node at (\n,.13) {\(\n\)};
    \node[text=unipd-muted] at (\n,-.16) {\(\times\)};
    \draw[unipd-primary] ({\n-.38},-2) rectangle ({\n+.38},-1.3);
    \node at (\n,-1.52) {\(\n\)};
    \draw[->,unipd-muted] (\n,-.4) -- ({\n+1},-1.25);
  }
  \foreach \n in {1,...,5} {
    \node[text=unipd-muted] at (\n,-1.81) {\(\times\)};
  }
  \node at (6,0) {\(\dots\)};
  \node at (6,-1.65) {\(\dots\)};
  \node[right] at (6.2,-.8) {\(n\mapsto n+1\)};
  \node[below] at (0,-2.05) {libera};
\end{tikzpicture}

  \caption{Schema dell'albergo di Hilbert: dopo lo
  spostamento \(n\mapsto n+1\), la camera \(0\)
  risulta libera. Le crocette indicano le camere
  occupate.}
  \label{fig:lez07-hilbert}
\end{figure}
\end{example}


\section{Operazioni sulle funzioni}
\label{sec:lez07-operazioni}

Si considerino due funzioni a valori reali
\[
  \psi\colon A\longrightarrow\R,\qquad
  \varphi\colon B\longrightarrow\R,
  \qquad A,B\subseteq\R.
\]
Per la \cref{def:lez06-dominio-codominio}, i domini
di \(\psi\) e \(\varphi\) sono rispettivamente \(A\)
e \(B\). La somma e il prodotto dei loro valori
possono quindi essere calcolati per ogni
\(x\in A\cap B\).

Le funzioni così ottenute hanno codominio
\(\R\), mentre la loro immagine, secondo la
\cref{def:lez06-immagine}, è costituita dai
valori effettivamente assunti.

\begin{definition}[Somma di funzioni]
\label{def:lez07-somma}
La \emph{somma} di \(\psi\) e \(\varphi\)
è la funzione
\[
  \psi+\varphi\colon A\cap B\longrightarrow\R,
  \qquad
  (\psi+\varphi)(x):=\psi(x)+\varphi(x).
\]
In particolare,
\(D(\psi+\varphi)=D(\psi)\cap D(\varphi)\).
\end{definition}

La somma è definita punto per punto mediante
l'addizione in \(\R\). La proprietà commutativa,
già introdotta per \(\N\) nella
\cref{sec:lez01-numeri-naturali}, vale anche in
\(\R\). Poiché le funzioni \(\psi+\varphi\) e
\(\varphi+\psi\) hanno lo stesso dominio e
codominio, per ogni \(x\in A\cap B\) si ha
\[
  (\psi+\varphi)(x)
  =\psi(x)+\varphi(x)
  =\varphi(x)+\psi(x)
  =(\varphi+\psi)(x).
\]
Pertanto \(\psi+\varphi=\varphi+\psi\).

\begin{example}[Somma di una radice e di una funzione lineare]
\label{ex:lez07-somma-radice}
Siano \(\varphi\colon[0,+\infty)\to\R\),
\(\varphi(x)=\sqrt{x}\), e
\(\psi\colon\R\to\R\), \(\psi(x)=2x\). Si ha
\[
  D(\psi+\varphi)
  =\R\cap[0,+\infty)
  =[0,+\infty),
\]
e, per ogni \(x\ge0\),
\[
  (\psi+\varphi)(x)=2x+\sqrt{x}.
\]
La somma non è definita per \(x<0\), poiché
tali valori non appartengono al dominio
di \(\varphi\).
\end{example}

\begin{definition}[Prodotto di funzioni]
\label{def:lez07-prodotto}
Il \emph{prodotto} di \(\psi\) e \(\varphi\)
è la funzione
\[
  \psi\varphi\colon A\cap B\longrightarrow\R,
  \qquad
  (\psi\varphi)(x):=\psi(x)\varphi(x).
\]
In particolare,
\(D(\psi\varphi)=D(\psi)\cap D(\varphi)\).
\end{definition}

Il prodotto è definito punto per punto mediante
la moltiplicazione in \(\R\), che è commutativa.
Come per la somma, le funzioni \(\psi\varphi\) e
\(\varphi\psi\) hanno lo stesso dominio e
codominio e, per ogni \(x\in A\cap B\), si ha
\[
  (\psi\varphi)(x)
  =\psi(x)\varphi(x)
  =\varphi(x)\psi(x)
  =(\varphi\psi)(x).
\]
Pertanto \(\psi\varphi=\varphi\psi\).

\begin{definition}[Rapporto di funzioni]
\label{def:lez07-quoziente}
Il \emph{rapporto} di \(\psi\) e \(\varphi\)
è la funzione
\[
  \frac{\psi}{\varphi}\colon
  (A\cap B)\setminus Z_\varphi
  \longrightarrow\R,
  \qquad
  \left(\frac{\psi}{\varphi}\right)(x)
  :=\frac{\psi(x)}{\varphi(x)},
\]
dove
\[
  Z_\varphi:=\{x\in D(\varphi):\varphi(x)=0\}
\]
è l'insieme degli zeri di \(\varphi\).
In particolare,
\[
  D(\psi/\varphi)
  =(D(\psi)\cap D(\varphi))\setminus Z_\varphi.
\]
\end{definition}

Il denominatore deve essere diverso da zero:
per questo gli zeri di \(\varphi\) sono esclusi
dal dominio comune.

Scambiando numeratore e denominatore si ottiene
\[
  D(\varphi/\psi)
  =(D(\psi)\cap D(\varphi))\setminus Z_\psi,
\]
dove \(Z_\psi\) indica l'insieme degli zeri
di \(\psi\).

A differenza della somma e del prodotto, il rapporto
non è commutativo: i due domini possono essere
diversi e, anche nei punti in cui entrambi i
rapporti sono definiti, i loro valori non
coincidono in generale.

\section{Grafici e trasformazioni delle funzioni}
\label{sec:lez07-trasformazioni}

Nella \cref{sec:lez06-applicazioni}, il grafico
di un'applicazione è stato introdotto come insieme
delle coppie ordinate che la rappresentano.
Si considera ora il caso delle funzioni reali
di variabile reale.

\begin{definition}[Grafico di una funzione]
\label{def:lez07-grafico}
Sia \(\varphi\colon A\to\R\), con \(A\subseteq\R\).
Il \emph{grafico} di \(\varphi\) è l'insieme
\[
  G(\varphi):=\{(x,\varphi(x)):x\in A\}
  \subseteq\R^2.
\]
\end{definition}

Ogni punto del grafico ha ascissa \(x\in A\) e
ordinata \(\varphi(x)\). Le trasformazioni considerate
nel seguito agiscono sulle ascisse o sulle ordinate,
permettendo di ottenere nuovi grafici a partire da
\(G(\varphi)\).

\subsection{Traslazione rigida}
\label{sec:lez07-traslazioni-verticali}

Sia \(\varphi\colon A\longrightarrow\R\), con
\(A\subseteq\R\), e sia \(c\in\R\). Si consideri
la funzione
\[
  \varphi+c\colon A\longrightarrow\R,
  \qquad
  (\varphi+c)(x):=\varphi(x)+c.
\]
Poiché la costante si somma ai valori di
\(\varphi\), il dominio rimane invariato:
\[
  D(\varphi+c)=D(\varphi)=A.
\]

Per la \cref{def:lez07-grafico}, il grafico
della nuova funzione è
\[
  G(\varphi+c)
  =\{(x,\varphi(x)+c):x\in A\}.
\]
A ogni punto \((x,\varphi(x))\in G(\varphi)\)
corrisponde quindi il punto
\[
  (x,\varphi(x)+c)=(x,\varphi(x))+(0,c).
\]
Viceversa, ogni punto di \(G(\varphi+c)\)
si ottiene in questo modo da un punto di
\(G(\varphi)\). Il nuovo grafico è pertanto
ottenuto mediante una \emph{traslazione rigida}
verticale secondo il vettore \((0,c)\).

Se \(c>0\), si ha
\[
  \varphi(x)+c>\varphi(x),
\]
e il grafico viene traslato verso l'alto.
Se \(c<0\), la disuguaglianza si inverte
e la traslazione avviene verso il basso.

In entrambi i casi, la distanza verticale
tra i punti corrispondenti è costante:
\[
  |(\varphi+c)(x)-\varphi(x)|=|c|.
\]
Per \(c=0\), la funzione e il suo grafico
rimangono invariati.

\subsection{Riscalamento verticale e ribaltamento}
\label{sec:lez07-riscalamento}

Sia \(\varphi\colon A\longrightarrow\R\), con
\(A\subseteq\R\), e sia \(b\in\R\). Si consideri
la funzione
\[
  b\varphi\colon A\longrightarrow\R,
  \qquad
  (b\varphi)(x):=b\,\varphi(x).
\]
Poiché il coefficiente \(b\) moltiplica i valori
di \(\varphi\), il dominio rimane invariato:
\[
  D(b\varphi)=D(\varphi)=A.
\]

Per la \cref{def:lez07-grafico}, il grafico
della nuova funzione è
\[
  G(b\varphi)=\{(x,b\varphi(x)):x\in A\}.
\]
La corrispondenza tra i due grafici è quindi
\[
  (x,\varphi(x))\longmapsto(x,b\varphi(x)).
\]
L'ascissa rimane invariata, mentre l'ordinata
viene moltiplicata per \(b\).

La distanza di un punto dall'asse delle ascisse
passa da \(|\varphi(x)|\) a
\[
  |b\varphi(x)|=|b|\,|\varphi(x)|.
\]
Ne segue che, se \(|b|>1\), il grafico subisce
una dilatazione verticale; se \(0<|b|<1\), si ha
invece una compressione verticale. I punti di
ordinata nulla rimangono fissi.

Se \(b<0\), oltre al riscalamento si ha un
cambiamento di segno delle ordinate. In
particolare, per \(b=-1\) si ottiene la funzione
\[
  -\varphi\colon A\longrightarrow\R,
  \qquad
  (-\varphi)(x)=-\varphi(x),
\]
e il grafico corrispondente si ottiene mediante
un \emph{ribaltamento} rispetto all'asse delle
ascisse.

Per \(b=1\), il grafico rimane invariato. Per
\(b=0\), si ottiene invece
\[
  G(0\varphi)=\{(x,0):x\in A\},
\]
ossia la parte dell'asse delle ascisse le cui
ascisse appartengono al dominio di \(\varphi\).

Non si può dedurre da \(b>1\) che
\(b\varphi(x)>\varphi(x)\) per ogni \(x\), poiché
il confronto dipende dal segno di \(\varphi(x)\).
La distanza verticale tra i punti corrispondenti
dei grafici di \(\varphi\) e \(b\varphi\) è
\[
  |b\varphi(x)-\varphi(x)|
  =|b-1|\,|\varphi(x)|.
\]

\begin{figure}[htbp]
  \centering
  \begin{tikzpicture}[unipd/diagram,x=1.2cm,y=1.2cm]
  \begin{scope}
    \draw[->,unipd-muted] (-1.4,0) -- (1.65,0) node[right] {\(x\)};
    \draw[->,unipd-muted] (0,-.3) -- (0,2.5) node[above] {\(y\)};
    \draw[very thick,unipd-primary,domain=-1:1,samples=50]
      plot (\x,{\x*\x});
    \draw[very thick,unipd-accent,domain=-1:1,samples=50]
      plot (\x,{\x*\x+1});
    \draw[->,unipd-muted,dashed] (.65,.4225) -- (.65,1.4225);
    \node[right=5pt,text=unipd-primary] at (1.02,1.08) {\(\varphi\)};
    \node[right=5pt,text=unipd-accent] at (1.02,2.08) {\(\varphi+1\)};
    \draw[unipd-muted] (.05,1) -- (-.05,1);
    \node[below left=3pt] at (0,1) {\(1\)};
  \end{scope}
  \begin{scope}[xshift=6.5cm]
    \draw[->,unipd-muted] (-1.4,0) -- (1.65,0) node[right] {\(x\)};
    \draw[->,unipd-muted] (0,-1.4) -- (0,1.5) node[above] {\(y\)};
    \draw[very thick,unipd-primary,domain=-1:1,samples=50]
      plot (\x,{\x*\x});
    \draw[very thick,unipd-accent,domain=-1:1,samples=50]
      plot (\x,{-\x*\x});
    \node[right=5pt,text=unipd-primary] at (1.02,1.08) {\(\varphi\)};
    \node[right=5pt,text=unipd-accent] at (1.02,-1.08) {\(-\varphi\)};
    \draw[unipd-muted] (.05,1) -- (-.05,1) node[left] {\(1\)};
    \draw[unipd-muted] (.05,-1) -- (-.05,-1) node[left] {\(-1\)};
  \end{scope}
\end{tikzpicture}

  \caption{A sinistra, traslazione rigida del
  grafico di \(\varphi(x)=x^2\) verso l'alto di
  un'unità; a destra, ribaltamento rispetto
  all'asse delle ascisse, da \(\varphi(x)=x^2\)
  a \(-\varphi(x)=-x^2\). Le curve sono
  rappresentate per \(x\in[-1,1]\).}
  \label{fig:lez07-verticali}
\end{figure}

\subsection{Traslazione orizzontale}
\label{sec:lez07-traslazioni-orizzontali}

Sia \(\varphi\colon A\longrightarrow\R\), con
\(A\subseteq\R\), e sia \(c\in\R\). Si consideri
la funzione
\[
  \psi\colon A_c\longrightarrow\R,
  \qquad
  \psi(x):=\varphi(x+c),
\]
dove
\[
  A_c:=\{x\in\R:x+c\in A\}
      =\{t-c:t\in A\}.
\]
Poiché l'argomento di \(\varphi\) deve appartenere
ad \(A\), il dominio della nuova funzione è
\[
  D(\psi)=A_c.
\]

Per la \cref{def:lez07-grafico}, il grafico
della nuova funzione è
\[
  G(\psi)=\{(t-c,\varphi(t)):t\in A\}.
\]
A ogni punto \((t,\varphi(t))\in G(\varphi)\)
corrisponde quindi il punto
\[
  (t-c,\varphi(t))
  =(t,\varphi(t))+(-c,0).
\]
Viceversa, ogni punto di \(G(\psi)\) si ottiene
in questo modo da un punto di \(G(\varphi)\).
Il nuovo grafico è pertanto ottenuto mediante
una \emph{traslazione rigida} orizzontale secondo
il vettore \((-c,0)\).

Se \(c>0\), il grafico viene traslato verso
sinistra; se \(c<0\), la traslazione avviene
verso destra. In entrambi i casi, la distanza
orizzontale tra i punti corrispondenti è
\[
  |(t-c)-t|=|c|.
\]
Per \(c=0\), la funzione e il suo grafico
rimangono invariati.

\begin{example}[Traslazione del dominio di un intervallo]
\label{ex:lez07-dominio-traslato}
Sia \(\varphi\colon[-1,1]\longrightarrow\R\).
Si considerino le funzioni
\[
  \psi_1(x):=\varphi(x+2),
  \qquad
  \psi_2(x):=\varphi(x-2).
\]

Per \(\psi_1\), la condizione sull'argomento è
\[
  -1\le x+2\le1
  \quad\Longleftrightarrow\quad
  -3\le x\le-1,
\]
quindi
\[
  D(\psi_1)=[-3,-1].
\]

Analogamente, per \(\psi_2\) si ha
\[
  -1\le x-2\le1
  \quad\Longleftrightarrow\quad
  1\le x\le3,
\]
da cui
\[
  D(\psi_2)=[1,3].
\]
\end{example}

\begin{figure}[htbp]
  \centering
  \begin{tikzpicture}[unipd/diagram,x=1.55cm,y=1.3cm]
  \draw[->,unipd-muted] (-3.45,0) -- (3.55,0) node[right] {\(x\)};
  \draw[->,unipd-muted] (0,-.4) -- (0,1.8) node[above] {\(y\)};
  \draw[very thick,unipd-primary,domain=-1:1,samples=50]
    plot (\x,{\x*\x});
  \draw[very thick,unipd-accent,domain=-3:-1,samples=50]
    plot (\x,{(\x+2)*(\x+2)});
  \draw[very thick,unipd-accent,dashed,domain=1:3,samples=50]
    plot (\x,{(\x-2)*(\x-2)});
  \foreach \x in {-3,-1,1,3} {
    \fill[unipd-primary] (\x,1) circle (1.8pt);
  }
  \foreach \x in {-3,-2,-1,1,2,3} {
    \draw[unipd-muted] (\x,.05) -- (\x,-.05) node[below] {\(\x\)};
  }
  \node[above,text=unipd-accent] at (-2,1.25) {\(\varphi(x+2)\)};
  \node[above,text=unipd-primary] at (0.35,1.25) {\(\varphi(x)\)};
  \node[above,text=unipd-accent] at (2,1.25) {\(\varphi(x-2)\)};
\end{tikzpicture}

  \caption{Traslazioni orizzontali del grafico
  di \(\varphi(x)=x^2\), con dominio \([-1,1]\).
  La funzione \(\psi_1(x)=\varphi(x+2)\) ha dominio
  \([-3,-1]\), mentre \(\psi_2(x)=\varphi(x-2)\)
  ha dominio \([1,3]\). I punti pieni indicano
  gli estremi inclusi.}
  \label{fig:lez07-orizzontali}
\end{figure}

\section{Intorni e punti di accumulazione}
\label{sec:lez07-intorni-accumulazione}

Per studiare il comportamento di una funzione
in prossimità di un punto, si introducono le
nozioni di intorno e di punto di accumulazione,
che saranno utilizzate nella definizione di limite.

\subsection{Intorni di punti reali e dell'infinito}
\label{sec:lez07-intorni}

\begin{definition}[Intorno circolare]
\label{def:lez07-intorno}
Siano \(x_0\in\R\) e \(\delta>0\).
L'\emph{intorno circolare} di centro \(x_0\)
e raggio \(\delta\) è l'insieme
\[
  D(x_0,\delta)
  :=\{x\in\R:|x-x_0|<\delta\}
  =(x_0-\delta,x_0+\delta).
\]
\end{definition}

La condizione \(|x-x_0|<\delta\) individua
i punti della retta reale la cui distanza
da \(x_0\) è inferiore a \(\delta\).
Il centro appartiene all'intorno, mentre
gli estremi \(x_0-\delta\) e \(x_0+\delta\)
sono esclusi.

\begin{figure}[htbp]
  \centering
  \begin{tikzpicture}[unipd/diagram,x=1.6cm]
  \draw[->,unipd-muted] (-3,0) -- (3.2,0) node[right] {\(x\)};
  \draw[very thick,unipd-primary] (-2,0) -- (2,0);
  \foreach \x in {-2,2} {
    \draw[unipd-primary,fill=white,thick] (\x,0) circle (2.4pt);
  }
  \fill[unipd-primary] (0,0) circle (1.8pt);
  \node[below=5pt] at (-2,0) {\(x_0-\delta\)};
  \node[below=5pt] at (0,0) {\(x_0\)};
  \node[below=5pt] at (2,0) {\(x_0+\delta\)};
  \draw[<->,unipd-muted] (-2,.5) -- (0,.5) node[midway,above] {\(\delta\)};
  \draw[<->,unipd-muted] (0,.5) -- (2,.5) node[midway,above] {\(\delta\)};
\end{tikzpicture}

  \caption{Intorno circolare di centro \(x_0\)
  e raggio \(\delta>0\).}
  \label{fig:lez07-intorno}
\end{figure}

\begin{definition}[Intorni di infinito]
\label{def:lez07-intorni-infinito}
Si dice \emph{intorno di \(+\infty\)}
ogni intervallo della forma
\[
  (a,+\infty)=\{x\in\R:x>a\},
  \qquad a\in\R.
\]
Analogamente, si dice \emph{intorno di \(-\infty\)}
ogni intervallo della forma
\[
  (-\infty,a)=\{x\in\R:x<a\},
  \qquad a\in\R.
\]
\end{definition}

A differenza degli intorni circolari,
gli intorni di infinito sono intervalli
illimitati della retta reale. I simboli
\(+\infty\) e \(-\infty\) non rappresentano
numeri reali e non appartengono agli
intervalli così definiti.

\begin{definition}[Retta reale estesa]
\label{def:lez07-retta-estesa}
La \emph{retta reale estesa} è l'insieme
\[
  \overline{\R}:=\R\cup\{-\infty,+\infty\},
\]
ordinato ponendo
\[
  -\infty<x<+\infty
  \qquad\forall x\in\R.
\]
\end{definition}

\begin{example}[Estremi di intervalli illimitati]
\label{ex:lez07-estremi-infiniti}
Per ogni \(a\in\R\), secondo le convenzioni
introdotte nella \cref{sec:lez04-estremi}, si ha
\[
  \sup(a,+\infty)=+\infty,
  \qquad
  \inf(a,+\infty)=a,
\]
\[
  \sup(-\infty,a)=a,
  \qquad
  \inf(-\infty,a)=-\infty.
\]

L'intervallo \((a,+\infty)\) non è
superiormente limitato, mentre
\((-\infty,a)\) non è inferiormente limitato.
Gli estremi infiniti appartengono alla
retta reale estesa, mentre gli estremi
finiti coincidono con \(a\), che non
appartiene a nessuno dei due intervalli.

Pertanto, nessuno dei due intervalli
ammette massimo o minimo.
\end{example}


\subsection{Punti di accumulazione}
\label{sec:lez07-punti-accumulazione}

\begin{definition}[Punto di accumulazione]
\label{def:lez07-accumulazione}
Sia \(A\subseteq\R\) non vuoto. Un punto
\(x_0\in\R\) si dice \emph{punto di accumulazione}
di \(A\) se
\[
  \forall\delta>0,\qquad
  \bigl(D(x_0,\delta)\cap A\bigr)\setminus\{x_0\}
  \neq\varnothing.
\]
L'insieme dei punti di accumulazione di \(A\)
si indica con \(A'\).
\end{definition}

Per la \cref{def:lez07-intorno}, la condizione
richiede che ogni intorno circolare di centro
\(x_0\) contenga almeno un elemento di \(A\)
distinto dal centro. Equivalentemente,
\[
  \forall\delta>0\ \exists a\in A:
  \qquad 0<|a-x_0|<\delta.
\]

La definizione non richiede che \(x_0\in A\).
Un punto può quindi essere di accumulazione
anche senza appartenere all'insieme.

Viceversa, un punto \(x_0\in A\) si dice
\emph{isolato} se esiste un intorno circolare
che non contiene altri elementi di \(A\).
In generale, negando la condizione della
\cref{def:lez07-accumulazione}, si ottiene
\[
  x_0\notin A'
  \quad\Longleftrightarrow\quad
  \exists\delta>0:
  \bigl(D(x_0,\delta)\cap A\bigr)
  \setminus\{x_0\}=\varnothing.
\]

Per dimostrare che un punto è di accumulazione
occorre dunque verificare la condizione per
ogni \(\delta>0\); per escluderlo è sufficiente
individuare un solo raggio positivo per cui
la condizione non sia soddisfatta.

\begin{example}[Un intervallo aperto e un punto isolato]
\label{ex:lez07-accumulazione-intervallo}
Sia
\[
  A=(1,2)\cup\{3\}.
\]
L'insieme dei suoi punti di accumulazione è
\[
  A'=[1,2].
\]
\end{example}

\begin{proof}
Si dimostra anzitutto che ogni punto di
\([1,2]\) è di accumulazione.

Sia \(x_0\in(1,2)\). Per ogni \(\delta>0\),
si ponga
\[
  \eta=
  \min\left\{
    \frac{\delta}{2},
    \frac{2-x_0}{2}
  \right\},
  \qquad
  a=x_0+\eta.
\]
Poiché \(\eta>0\), si ha \(a\in(1,2)\)
e
\[
  0<|a-x_0|=\eta<\delta.
\]
Per la \cref{def:lez07-accumulazione},
\(x_0\in A'\).

Anche gli estremi \(1\) e \(2\) sono punti
di accumulazione. Per ogni \(\delta>0\),
posto
\[
  \eta=\min\left\{\frac{\delta}{2},
                  \frac12\right\},
\]
i punti \(1+\eta\) e \(2-\eta\) appartengono
ad \(A\) e distano rispettivamente da
\(1\) e da \(2\) una quantità positiva
minore di \(\delta\).

Pertanto
\[
  [1,2]\subseteq A'.
\]

Si dimostra ora che nessun punto esterno
a \([1,2]\) è di accumulazione.

Se \(x_0<1\), scegliendo
\[
  \delta=\frac{1-x_0}{2}>0
\]
si ottiene
\[
  D(x_0,\delta)\cap A=\varnothing.
\]
Quindi \(x_0\notin A'\).

Se \(x_0>2\) e \(x_0\ne3\), si scelga
\[
  \delta=
  \frac12\min\{x_0-2,|x_0-3|\}>0.
\]
L'intorno \(D(x_0,\delta)\) non interseca
l'intervallo \((1,2)\) e non contiene
il punto \(3\). Pertanto
\[
  D(x_0,\delta)\cap A=\varnothing,
\]
da cui \(x_0\notin A'\).

Infine, per \(x_0=3\), scegliendo
\(\delta=1/2\) si ha
\[
  D(3,1/2)\cap A=\{3\},
\]
quindi
\[
  \bigl(D(3,1/2)\cap A\bigr)
  \setminus\{3\}=\varnothing.
\]
Il punto \(3\) è dunque isolato e non
è di accumulazione.

Si è così dimostrato che
\(A'\subseteq[1,2]\). Dalla doppia
inclusione segue
\[
  A'=[1,2].
\]
\end{proof}

\begin{figure}[htbp]
  \centering
  \begin{tikzpicture}[unipd/diagram,x=2.4cm,y=1cm]
  \foreach \y in {0,-1.3} {
    \draw[->,unipd-muted] (.55,\y) -- (3.5,\y) node[right] {\(x\)};
    \foreach \x in {1,2,3} {
      \draw[unipd-muted] (\x,{\y+.06}) -- (\x,{\y-.06})
        node[below=4pt] {\(\x\)};
    }
  }
  \node[left] at (.45,0) {\(A\)};
  \node[left] at (.45,-1.3) {\(A'\)};
  \draw[very thick,unipd-primary] (1,0) -- (2,0);
  \foreach \x in {1,2} {
    \draw[unipd-primary,fill=white,thick] (\x,0) circle (2.4pt);
  }
  \fill[unipd-primary] (3,0) circle (2.4pt);
  \draw[very thick,unipd-accent] (1,-1.3) -- (2,-1.3);
  \foreach \x in {1,2} {
    \fill[unipd-accent] (\x,-1.3) circle (2.4pt);
  }
\end{tikzpicture}

  \caption{L'insieme \(A=(1,2)\cup\{3\}\)
  e il suo insieme dei punti di accumulazione
  \(A'=[1,2]\). I cerchi vuoti indicano
  estremi esclusi, mentre i punti pieni
  indicano elementi inclusi.}
  \label{fig:lez07-accumulazione}
\end{figure}

\begin{exercise}[Verifica dei punti di accumulazione]
\label{ex:lez07-verifica-accumulazione}
Per \(A=(1,2)\cup\{3\}\), verificare direttamente
che \(2\) e ogni \(x_0\in(1,2)\) sono punti
di accumulazione, mentre ogni \(x_0>2\),
con \(x_0\ne3\), non lo è.
\end{exercise}

\begin{proof}[Soluzione dell'esercizio]
Nella dimostrazione dell'\cref{ex:lez07-accumulazione-intervallo},
per ogni \(x_0\in(1,2]\) e ogni \(\delta>0\)
è stato costruito un elemento \(a\in A\) tale che
\(0<|a-x_0|<\delta\). Per \(x_0>2\), con
\(x_0\ne3\), è stato invece individuato un
intorno disgiunto da \(A\). Le conclusioni
seguono quindi dalla \cref{def:lez07-accumulazione}.
\end{proof}

\begin{exercise}[Variante con estremi inclusi]
\label{ex:lez07-accumulazione-chiuso}
Ripetere la verifica dell'\cref{ex:lez07-verifica-accumulazione}
per l'insieme
\[
  B=[1,2]\cup\{3\}
\]
e determinarne tutti i punti di accumulazione.
\end{exercise}

\begin{proof}[Soluzione dell'esercizio]
Si ha ancora
\[
  B'=[1,2].
\]

Infatti, poiché \((1,2)\subseteq B\),
le costruzioni utilizzate nella dimostrazione
dell'\cref{ex:lez07-accumulazione-intervallo}
rimangono valide per ogni \(x_0\in[1,2]\).
Anche per gli estremi \(1\) e \(2\), gli elementi
scelti sono distinti dal centro dell'intorno.

Se \(x_0<1\) oppure \(x_0>2\), con \(x_0\ne3\),
gli stessi raggi consentono di individuare
un intorno disgiunto da \(B\).

Infine,
\[
  D(3,1/2)\cap B=\{3\},
\]
quindi \(3\) è ancora un punto isolato.

L'inclusione degli estremi nell'intervallo
modifica la loro appartenenza all'insieme,
ma non la loro proprietà di essere punti
di accumulazione.
\end{proof}
